В выводе мультимножество записывается в виде \code{⟨2(000), 1(011)⟩}: перед словом в скобках стоит его показатель. Слова с нулевыми показателями пропускаются, пустое мультимножество обозначается \code{⟨⟩}. При ручном вводе задаются пары \code{слово:показатель} через пробел, например \code{000:2 011:1}; если показатель не указан, он равен 1.

Носитель хранится в классе \code{CodeSet} (подраздел~\ref{sec:codeset}). Поле \code{depth} задаёт разрядность $n$. Контейнер \code{codes} типа \code{std::vector<size_t>} содержит числовые значения слов в порядке кода Грея; двоичная запись нужной длины формируется при выводе слова. Вектор \code{positions} того же типа задаёт обратное соответствие: по числовому значению слова определяется его номер в \code{codes}.

Мультимножество реализовано классом \code{Multiset} (подраздел~\ref{sec:multiset}). Поле \code{carrier} типа \code{std::shared_ptr<const CodeSet>} задаёт общий носитель для $\hat U$, $A$, $B$ и результатов операций. Слова не копируются в каждое мультимножество, а изменение общего носителя через \code{carrier} запрещено. В поле \code{counts} типа \code{std::vector<size_t>} каждое мультимножество хранит собственные показатели всех $N$ слов, включая нулевые.

В описаниях методов \code{this} — текущий экземпляр соответствующего класса: для \code{CodeSet} это носитель, для \code{Multiset} — мультимножество. Индекс $j$ в векторах начинается с 0 и соответствует элементу $x_{j+1}$ в записи математического описания, начинающейся с $x_1$.

\subsection{Код Грея и носитель}
\label{sec:codeset}

\fn{changed_bit}
{номер шага $i>0$.}
{Номер разряда, который нужно инвертировать: число нулей в конце двоичной записи $i$.}
\functioncode{changed_bit}{01.cpp}

\fn{generate_gray_code}
{\code{depth} — разрядность $n$.}
{Последовательность $N=2^n$ слов кода Грея. Первое слово состоит из нулей; каждое следующее получается инверсией разряда, найденного \code{changed_bit}. При $n=0$ последовательность содержит одно пустое слово.}
\functioncode{generate_gray_code}{02.cpp}

\fn{code_to_string}
{\code{code} — слово кода; \code{depth} — его разрядность $n$.}
{Двоичная запись слова длины $n$, включая начальные нули.}
\functioncode{code_to_string}{03.cpp}

\fn{CodeSet::CodeSet}
{\code{this} — создаваемый носитель; \code{depth} — разрядность $n$.}
{Универсум $U$ из $N=2^n$ слов в порядке кода Грея и соответствие между каждым словом и его номером.}
\functioncode{CodeSet::CodeSet}{04.cpp}

\fn{CodeSet::get_depth}
{\code{this} — носитель $U$.}
{Разрядность $n$ слов носителя.}
\functioncode{CodeSet::get_depth}{05.cpp}

\fn{CodeSet::size}
{\code{this} — носитель $U$.}
{Число слов $N=2^n$.}
\functioncode{CodeSet::size}{06.cpp}

\fn{CodeSet::get_code}
{\code{this} — носитель $U$; \code{i} — номер слова, $0\le i<N$.}
{Слово $x_{i+1}$ длины $n$. При недопустимом номере возникает ошибка.}
\functioncode{CodeSet::get_code}{07.cpp}

\fn{CodeSet::find}
{\code{this} — носитель $U$; \code{code} — искомое слово.}
{Номер слова в порядке кода Грея либо признак отсутствия подходящего слова. Допустимое слово состоит из $n$ символов 0 и 1; при $n=0$ допустимо пустое слово.}
\functioncode{CodeSet::find}{08.cpp}

\subsection{Автоматический выбор показателей}

\fn{generator}
{текущее состояние генератора случайных чисел.}
{Генератор последовательности, используемой при выборе $k$, мощности мультимножества и слова, показатель которого увеличивается при автоматическом заполнении.}
\functioncode{generator}{09.cpp}

\fn{make_seed}
{значение источника случайных чисел и текущее время.}
{Начальный параметр $s$ генератора, применяемого для автоматического заполнения мультимножеств.}
\functioncode{make_seed}{10.cpp}

\fn{seed_random}
{\code{seed} — начальный параметр $s$.}
{Установлено начальное состояние последовательности случайных чисел для автоматического выбора показателей.}
\functioncode{seed_random}{11.cpp}

\fn{random_below}
{\code{bound} — положительная верхняя граница $b$.}
{Равновероятно выбранное целое число $r$, $0\le r<b$. Для показателя $k$ используется $1+r$ при $b=100$; при заполнении выбирается номер одного из ещё не заполненных слов.}
\functioncode{random_below}{12.cpp}

\subsection{Мультимножества и операции}
\label{sec:multiset}

Для методов этого подраздела \code{this} обозначает текущее мультимножество $\hat X$. При выполнении операций его роль играет $A$, а аргумент \code{b} обозначает $B$. Мультимножества должны иметь общий носитель. В формулах ниже $i=1,\ldots,N$.

\fn{Multiset::Multiset}
{\code{this} — создаваемое мультимножество; \code{carrier} — носитель $U$.}
{Пустое мультимножество $\hat X$ над $U$: $a_i(x_i)=0$ для всех слов, $m_{\hat X}=0$.}
\functioncode{Multiset::Multiset}{13.cpp}

\fn{Multiset::check_same_carrier}
{\code{this} — $\hat X$; \code{other} — $\hat Y$.}
{Проверено, что $\hat X$ и $\hat Y$ заданы над одним носителем $U$. Иначе возникает ошибка несовпадения носителей.}
\functioncode{Multiset::check_same_carrier}{14.cpp}

\fn{Multiset::check_within}
{\code{this} — $\hat X$; \code{universe} — $\hat U$.}
{Проверено включение $\hat X\subseteq\hat U$, то есть $a_i(x_i)\le a_{\hat U}(x_i)$ для каждого слова. При нарушении условия возникает ошибка.}
\functioncode{Multiset::check_within}{15.cpp}

\fn{Multiset::get_carrier}
{\code{this} — $\hat X$.}
{Носитель $U$, задающий слова и их порядок в $\hat X$.}
\functioncode{Multiset::get_carrier}{16.cpp}

\fn{Multiset::size}
{\code{this} — $\hat X$.}
{Число слов носителя $N$; оно не зависит от показателей и мощности $\hat X$.}
\functioncode{Multiset::size}{17.cpp}

\fn{Multiset::get_count}
{\code{this} — $\hat X$; \code{i} — номер слова, $0\le i<N$.}
{Показатель $a_{i+1}(x_{i+1})$. При недопустимом номере возникает ошибка.}
\functioncode{Multiset::get_count}{18.cpp}

\fn{Multiset::set_count}
{\code{this} — $\hat X$; \code{i} — номер слова; \code{count} — новый показатель $a\ge0$.}
{Показатель слова $x_{i+1}$ заменён на $a$. Допустимость номера проверяется здесь, а ограничение показателя значением в $\hat U$ — при заполнении мультимножества.}
\functioncode{Multiset::set_count}{19.cpp}

\fn{Multiset::cardinality}
{\code{this} — $\hat X$.}
{Мощность $m_{\hat X}=\sum_{i=1}^{N}a_i(x_i)$.}
\functioncode{Multiset::cardinality}{20.cpp}

\fn{Multiset::is_submultiset_of}
{\code{this} — $\hat X$; \code{other} — $\hat Y$.}
{Истинность включения $\hat X\subseteq\hat Y$: показатель каждого слова в $\hat X$ не больше показателя того же слова в $\hat Y$.}
\functioncode{Multiset::is_submultiset_of}{21.cpp}

\fn{Multiset::unite}
{\code{this} — $A$; \code{b} — $B$.}
{Объединение $A\cup B$ с показателями $\max\{\aA,\aB\}$.}
\functioncode{Multiset::unite}{22.cpp}

\fn{Multiset::intersect}
{\code{this} — $A$; \code{b} — $B$.}
{Пересечение $A\cap B$ с показателями $\min\{\aA,\aB\}$.}
\functioncode{Multiset::intersect}{23.cpp}

\fn{Multiset::complement}
{\code{this} — $A$; \code{universe} — $\hat U$.}
{Дополнение $\overline A$ с показателями $\aU-\aA$. Перед вычислением проверяется $A\subseteq\hat U$.}
\functioncode{Multiset::complement}{24.cpp}

\fn{Multiset::difference}
{\code{this} — $A$; \code{b} — $B$; \code{universe} — $\hat U$.}
{Разность $A\setminus B=A\cap\overline B$ с показателями $\min\{\aA,\aU-\aB\}$. При перестановке $A$ и $B$ получается $B\setminus A$.}
\functioncode{Multiset::difference}{25.cpp}

\fn{Multiset::symmetric_difference}
{\code{this} — $A$; \code{b} — $B$; \code{universe} — $\hat U$.}
{Симметрическая разность $(A\setminus B)\cup(B\setminus A)$. Показатель каждого слова равен большему из его показателей в двух разностях.}
\functioncode{Multiset::symmetric_difference}{26.cpp}

\fn{Multiset::arithmetic_sum}
{\code{this} — $A$; \code{b} — $B$; \code{universe} — $\hat U$.}
{Арифметическая сумма $A+B$ с показателями $\min\{\aA+\aB,\aU\}$. Проверяется включение обоих мультимножеств в $\hat U$.}
\functioncode{Multiset::arithmetic_sum}{27.cpp}

\fn{Multiset::arithmetic_difference}
{\code{this} — $A$; \code{b} — $B$.}
{Арифметическая разность $A-B$ с показателями $\max\{\aA-\aB,0\}$.}
\functioncode{Multiset::arithmetic_difference}{28.cpp}

\fn{Multiset::arithmetic_product}
{\code{this} — $A$; \code{b} — $B$; \code{universe} — $\hat U$.}
{Арифметическое произведение $A*B$ с показателями $\min\{\aA*\aB,\aU\}$. Проверяется включение обоих мультимножеств в $\hat U$.}
\functioncode{Multiset::arithmetic_product}{29.cpp}

\fn{Multiset::arithmetic_division}
{\code{this} — $A$; \code{b} — $B$.}
{Арифметическое частное $A:B$: показатель равен $\lfloor\aA:\aB\rfloor$ при $\aB\ne0$ и нулю при $\aB=0$.}
\functioncode{Multiset::arithmetic_division}{30.cpp}

\fn{uniform_universe}
{\code{carrier} — универсум $U$; \code{multiplicity} — показатель $k$, $1\le k\le100$.}
{Мультимножество $\hat U$ с одинаковыми показателями $a_{\hat U}(x_i)=k$ и мощностью $m_{\hat U}=kN$. Недопустимое значение $k$ вызывает ошибку.}
\functioncode{uniform_universe}{31.cpp}

\fn{random_submultiset}
{\code{bound} — ограничивающее мультимножество $\hat U$; \code{cardinality} — требуемая мощность $c$, $0\le c\le m_{\hat U}$.}
{Мультимножество мощности $c$ с показателями, не превышающими показателей в $\hat U$. На каждом шаге случайно выбирается слово со свободным местом и его показатель увеличивается на 1. После достижения предела слово исключается из дальнейшего выбора.}
\functioncode{random_submultiset}{32.cpp}

\subsection{Ошибки и ограничения}

При ручном вводе проверяются длина и состав слова, запись показателя, свободное место у слова и оставшаяся мощность. Изменения принимаются только после проверки всей строки. Пустая строка отменяет заполнение. Повторное создание универсума требует заново задать $A$ и $B$.

\fn{GrayCodeError::GrayCodeError}
{\code{this} — создаваемая ошибка; \code{message} — её описание.}
{Обозначена ошибка при работе с кодом Грея.}
\functioncode{GrayCodeError::GrayCodeError}{33.cpp}

\fn{InvalidDepthError::InvalidDepthError}
{\code{this} — создаваемая ошибка.}
{Сообщено, что заданная разрядность не позволяет представить $N=2^n$ в программе.}
\functioncode{InvalidDepthError::InvalidDepthError}{34.cpp}

\fn{CodeIndexError::CodeIndexError}
{\code{this} — создаваемая ошибка.}
{Сообщено, что номер слова находится вне носителя.}
\functioncode{CodeIndexError::CodeIndexError}{35.cpp}

\fn{MultisetError::MultisetError}
{\code{this} — создаваемая ошибка; \code{message} — её описание.}
{Обозначена ошибка при работе с мультимножеством.}
\functioncode{MultisetError::MultisetError}{36.cpp}

\fn{MultisetIndexError::MultisetIndexError}
{\code{this} — создаваемая ошибка.}
{Сообщено, что для указанного номера нет элемента носителя.}
\functioncode{MultisetIndexError::MultisetIndexError}{37.cpp}

\fn{CarrierMismatchError::CarrierMismatchError}
{\code{this} — создаваемая ошибка.}
{Сообщено о попытке выполнить операцию над мультимножествами с разными носителями.}
\functioncode{CarrierMismatchError::CarrierMismatchError}{38.cpp}

\fn{NotSubmultisetError::NotSubmultisetError}
{\code{this} — создаваемая ошибка.}
{Сообщено о превышении показателя слова относительно ограничения в $\hat U$.}
\functioncode{NotSubmultisetError::NotSubmultisetError}{39.cpp}

\fn{CardinalityOverflowError::CardinalityOverflowError}
{\code{this} — создаваемая ошибка.}
{Сообщено, что требуемая мощность $c$ превышает мощность ограничивающего мультимножества.}
\functioncode{CardinalityOverflowError::CardinalityOverflowError}{40.cpp}

\fn{InvalidMultiplicityError::InvalidMultiplicityError}
{\code{this} — создаваемая ошибка.}
{Сообщено, что показатель $k$ не принадлежит отрезку $[1,100]$.}
\functioncode{InvalidMultiplicityError::InvalidMultiplicityError}{41.cpp}

\fn{InputClosedError::InputClosedError}
{\code{this} — создаваемая ошибка.}
{Зафиксировано завершение пользовательского ввода.}
\functioncode{InputClosedError::InputClosedError}{42.cpp}

Через меню задаются $0\le n\le19$, $1\le k\le100$, $0\le m_A,m_B\le m_{\hat U}$. Числа вводятся записью не длиннее девяти цифр. При $n=0$ универсум содержит единственное пустое слово $\varepsilon$; для его ручного ввода используется обозначение \code{ε}.

\needspace{12\baselineskip}
\subsection{Функция main}

При входе в \code{main} происходит инициализация параметров.
\functioncode{main}{main.cpp}

\subsection{Время формирования универсума}

Программа \code{bench} измеряет создание универсума $U$.
\functioncode{main (tools/bench.cpp)}{bench_main.cpp}

\functioncode{seconds}{seconds.cpp}

Для каждой разрядности от 10 до 19 выполнено 21 измерение. В таблице~\ref{tab:bench} приведена медиана.

\begin{table}[H]
\caption{Зависимость времени формирования универсума от разрядности}
\label{tab:bench}
\centering
\begin{singlespace}
\begin{tabular}{|c|c|}
\hline
Разрядность $n$ & Время формирования, мс \\\hline
10 & 0,0017 \\\hline
11 & 0,0031 \\\hline
12 & 0,0058 \\\hline
13 & 0,0117 \\\hline
14 & 0,0245 \\\hline
15 & 0,0491 \\\hline
16 & 0,1558 \\\hline
17 & 0,3026 \\\hline
18 & 0,5981 \\\hline
19 & 1,2608 \\\hline
\end{tabular}
\end{singlespace}
\end{table}
